% #1091 修复前后对比用 MWE。第 1、2 行是 issue 中的写法，其余三行是同一根因的其他表现。
% 每行左侧的灰色文字只是标签；第 6 行起是不加装饰的直接输入，作为间距的对照。
\documentclass{article}
\usepackage[paperwidth=12cm,paperheight=6.2cm,margin=4mm]{geometry}
\usepackage{xcolor}
\usepackage{xeCJK}
\usepackage{xeCJKfntef}
\pagestyle{empty}
\setlength\parindent{0pt}
\setlength\parskip{3pt}

\newsavebox{\tri}
\newcommand{\myfillin}{%
  \sbox{\tri}{\rule{0.8em}{0.5em}}
  \CJKunderline{\hspace*{0.5em plus .5em minus .5em}\usebox{\tri}\hspace*{0.5em plus .5em minus .5em}}}

\newcommand\Label[1]{\makebox[2.6cm][l]{\footnotesize\color{gray}#1}\kern0pt\relax}

\begin{document}
\Label{(1) 普通空格}普通字符 \myfillin{} 后续文字

\Label{(2) 使用 \texttt{\string~}}普通字符~\myfillin{} 后续文字

\Label{(3) 无空格填空线}今天是\CJKunderline{\hspace*{3em}}年

\Label{(4) 书名号开头}书\CJKunderline{《红楼梦》}的读后感

\Label{(5) 句号结尾}\CJKunderline{虚室生白。}吉祥止止

\Label{(6) 直接输入对照}书《红楼梦》的读后感

\Label{}虚室生白。吉祥止止
\end{document}
